\documentclass[11pt]{article}
\usepackage[margin=1.05in]{geometry}
\usepackage{amsmath,amsthm,amssymb,mathtools,bm}
\usepackage{booktabs,array,enumitem}\usepackage[expansion=false]{microtype}
\usepackage{xcolor}
\usepackage[colorlinks=true,linkcolor=blue!60!black,citecolor=blue!60!black,urlcolor=blue!60!black]{hyperref}
\theoremstyle{plain}
\newtheorem{theorem}{Theorem}[section]\newtheorem{proposition}[theorem]{Proposition}\newtheorem{lemma}[theorem]{Lemma}
\newtheorem{corollary}[theorem]{Corollary}\newtheorem{conjecture}[theorem]{Conjecture}
\theoremstyle{definition}\newtheorem{definition}[theorem]{Definition}\newtheorem{remark}[theorem]{Remark}
\newtheorem{example}[theorem]{Example}
\newcommand{\R}{\mathbb R}\newcommand{\Z}{\mathbb Z}\newcommand{\E}{\mathbb E}\newcommand{\Var}{\operatorname{Var}}\newcommand{\Cov}{\operatorname{Cov}}
\newcommand{\relu}{\operatorname{relu}}\newcommand{\diag}{\operatorname{diag}}\newcommand{\tr}{\operatorname{tr}}
\newcommand{\Tr}{\operatorname{Tr}}\newcommand{\1}{\mathbf 1}\newcommand{\cA}{\mathcal A}\newcommand{\cB}{\mathcal B}
\newcommand{\cF}{\mathcal F}\newcommand{\cT}{\mathcal T}\newcommand{\cH}{\mathcal H}\newcommand{\aff}{\mathfrak{aff}}
\newcommand{\He}{\mathrm{He}}\newcommand{\Fa}{\Phi_{\mathrm F}}\newcommand{\cZ}{\mathcal Z}\newcommand{\Sym}{\operatorname{Sym}}
\newcommand{\Aff}{\mathrm{Aff}}\newcommand{\cN}{\mathcal N}
\newcommand{\F}{\mathbb F}\newcommand{\Ad}{\operatorname{Ad}}\renewcommand{\arg}{\operatorname{arg}}
\title{\textsc{Localization in the commutant, stage 17:}\\ \large a study of ``An isomorphism of the free group factors''.\\
The argument, its five constructs, its essence,\\ and what it changes for the programme}
\author{}
\date{10 October 2026}
\begin{document}
\maketitle
\begin{abstract}
The paper (OpenAI, 23 September 2026) proves $L(\F_n)\cong L(\F_{n+1})$ for every $n\ge3$. With Dykema's amplification formula this gives
$L(\F_2)\cong L(\F_3)$, and with the Dykema--Rădulescu alternative it shows that all interpolated free group factors are isomorphic, with
fundamental group $\R_{>0}$. Free entropy dimensions are then not invariants of the generated von Neumann algebra.
\begin{itemize}[leftmargin=*]
\item \emph{Reading.} We reconstruct the proof in five moves and check each step's logic.
\item \emph{Checks.} We test every computable lemma: the Fox cocycle and the prefix identity exactly; the free-Poisson spectral limit, the
cutoff inversion and the free-sum norm bound in random-matrix models; and the mechanism of the key flow (letters barely move while a long word
rotates by the extra generator).
\item \emph{Verdict.} We find no gap.
\end{itemize}
The essence, in one sentence: inside $L(\F_{n+1})$ there are trace-preserving automorphisms fixing the extra generator $C$ that move each of the
other $n$ generators by an arbitrarily small amount in operator norm, while rotating one very long word in them by almost exactly $C$.
Iterated with error budgets chosen after each step, they produce a norm-convergent sequence of free Haar $n$-tuples whose limit swallows
$C$, and therefore generates everything. Four mechanisms make it work:
\begin{itemize}[leftmargin=*]
\item a rich family of explicit automorphisms, the Fox-cocycle flows;
\item coherent $\sqrt m$ amplification of tiny per-letter rotations along words with free prefixes;
\item freeness turning $\ell^2$ smallness into operator-norm smallness;
\item the asymmetry between distribution, which is closed under norm limits, and generation, which is not.
\end{itemize}
We end with what this changes for our programme. Its central cocycle is literally our stage-15 error cocycle, and its adaptive iteration is a
precise template for the instance-coupled dynamics the user has been asking for.

\emph{Addendum} (Section~\ref{sec:deductions}). We confirm the companion study's deductions and sharpen them:
\begin{itemize}[leftmargin=*]
\item a relative absorption $P*L(\Z)\cong P*L(\F_{k+1})$ fixing $P$;
\item weak nullness of the discarded completions, which form a free Haar unitary in $M^\omega$;
\item a hard-edge steering law, $\|T_mh-\delta_e\|\,\|h\|\ge(1+o(1))/(2\sqrt m)$, which forces $m=\Theta(\varepsilon^{-4})$ on the paper's route against the
universal $2/\varepsilon-1$;
\item the Fox flows as cocycle self-conjugacies of the free Bernoulli shift $M=(*_\Gamma L(\Z))\rtimes\Gamma$.
\end{itemize}
\end{abstract}
\tableofcontents
\section{Summary and verdict}
\label{sec:summary}
\paragraph{What is proved.}
\begin{itemize}[leftmargin=*]
\item For $n\ge3$, $L(\F_{n+1})$ contains a freely generating Haar $n$-tuple: $n$ unitaries whose nonidentity reduced words all have trace zero and which
generate the whole factor. So $L(\F_n)\cong L(\F_{n+1})$.
\item Amplifying $L(\F_3)\cong L(\F_5)$ by $\sqrt2$ gives $L(\F_2)\cong L(\F_3)$.
\item Hence all interpolated free group factors $L(\F_r)$, $1<r\le\infty$, are isomorphic, and their fundamental group is $\R_{>0}$.
\item The free entropy dimensions $\delta,\delta_0,\delta^*,\delta^\star$ take every integer value $n\ge2$ on finite generating tuples of $L(\F_2)$.
\end{itemize}
\paragraph{How.} The construction has five moves, each elementary given freeness (Section~\ref{sec:moves}).
\begin{enumerate}[leftmargin=*]
\item Reduce isomorphism to finding a freely generating Haar tuple.
\item Take the logarithm $S=\arg C$ of the extra generator and its conjugates $S_g=A_gSA_g^*$. They form a free family indexed by $\F_n$.
\item Build polynomial flows $A_j'=i\,s(h_j)A_j$ with $s(k)=\sum_gk(g)A_gSA_g^*$, where the coefficients form a Fox cocycle.
\begin{itemize}
\item The flows are trace-preserving, fix $C$, and move each letter by at most $3\pi\|h_j\|_{\ell^2}$.
\item They move a word $A_w$ to within $3\pi\|D_w-\delta_e\|_{\ell^2}$ of $e^{iS}A_w$.
\end{itemize}
\item Make $h$ tiny while $D_w\approx\delta_e$. With $m$ free prefixes, $D_w=T_mh$ where $T_mT_m^*/m$ tends to the free Poisson law, which has no atom at $0$.
Truncated inversion gives $\|h\|_{\ell^2}^2\approx2/(\pi m\sqrt d)$ with error $\approx(2/\pi)\sqrt d$.
\item Absorb and iterate. Compose with the basis change $C\mapsto CA_w$, so the moved $n$-tuple contains a word close to $C$. Repeat with budgets chosen
after each word is known. The $n$-tuples converge in norm, keep the free Haar law, and their limit generates.
\end{enumerate}
\paragraph{Verdict.} We checked every step's logic and every computable lemma, and found no gap (Section~\ref{sec:critical}). The weight-bearing
steps are the trace-preservation argument (a first variation at a free point plus real-analytic continuation) and the spectral inversion
(free Poisson without atom). Both hold. The external inputs are classical: the amplification formula and the alternative.

\paragraph{For us.} The Fox cocycle $D_{gb}=D_g+\lambda(g)D_b$ is the same object as our stage-15 error cocycle. The paper's analytic heart, prefix
transports by free unitaries adding like a circular element, is the mechanism behind our measured sum rule. Its adaptive iteration, where each
tolerance is fixed only after the previous step's output is seen, is a rigorous template for the weight-coupled, dynamically adapted estimator
we have been seeking (Section~\ref{sec:programme}).

\section{The statement and the classical reduction}
\label{sec:statement}
For a discrete group $G$, $L(G)=\lambda(G)''$ on $\ell^2(G)$ with trace $\tau(x)=\langle x\delta_e,\delta_e\rangle$. The free group factor problem, attributed to Kadison, asks
whether $L(\F_m)\cong L(\F_n)$ for $m\ne n$. Three classical facts frame it.
\begin{itemize}[leftmargin=*]
\item \emph{Interpolated factors and amplification} (Dykema, Rădulescu). There is a family $N_s=L(\F_s)$, $s>1$, with
\[
N_s^{\,t}\cong N_{1+(s-1)t^{-2}}\qquad(s>1,\ t>0),
\]
where $Q^t$ is the corner of $Q\otimes B(\ell^2)$ of trace $t$.
\item \emph{The alternative} (Dykema; Rădulescu with $s=\infty$). Either all $N_s$, $1<s\le\infty$, are isomorphic, or they are pairwise
non-isomorphic.
\item \emph{Haar tuples} (Lemma~2.2 of the paper). If unitaries $U_1,\dots,U_d$ generate $M$ and every nonidentity reduced word has trace zero, then
$\lambda(x_j)\mapsto U_j$ extends to a trace-preserving isomorphism $L(\F_d)\to M$. The proof is that the words form an orthonormal basis of $L^2(M)$.
\end{itemize}
So it is enough to find, inside $L(\F_{n+1})$, $n$ unitaries with free Haar word distribution that generate it.
\begin{center}
\begin{tabular}{lll}\toprule
step & isomorphism & source\\\midrule
Theorem 1.1 at $n=3,4$ & $N_3\cong N_4\cong N_5$ & the construction\\
amplify by $\sqrt2$ & $N_3^{\sqrt2}\cong N_5^{\sqrt2}$ & functoriality\\
amplification formula & $N_3^{\sqrt2}=N_{1+2/2}=N_2$, $N_5^{\sqrt2}=N_{1+4/2}=N_3$ & Dykema\\
conclusion & $L(\F_2)\cong L(\F_3)$ & \\
alternative & all $N_r$ isomorphic, $\mathcal F(N_r)=\R_{>0}$ & Dykema--Rădulescu\\\bottomrule
\end{tabular}
\end{center}
The construction needs $n\ge3$ because its prefix words use three letters, $a=x_1$ and $b_j=x_2^jx_3x_2^{-j}$. Rank $2$ comes only through amplification.

\section{The construction in five moves}
\label{sec:moves}
Fix $n\ge3$ and $M=L(\F_{n+1})$ with canonical generators $(A_1,\dots,A_n,C)$, a freely generating Haar $(n+1)$-tuple. Write $\Gamma=\F_n=\langle x_1,\dots,x_n\rangle$, and
$A_g$ for the word $g\in\Gamma$ evaluated in the first $n$ letters.

\subsection{Move 1: the bounded logarithm and its free conjugate family}
Let $S=\arg C\in W^*(C)$, with values in $(-\pi,\pi]$. Then $S=S^*$, $\|S\|\le\pi$, $\tau(S)=0$, $e^{iS}=C$ and $W^*(S)=W^*(C)$. For $g\in\Gamma$ put $S_g=A_gSA_g^*$.
\begin{itemize}[leftmargin=*]
\item \emph{The conjugates are free} (Lemma 2.3). The algebras $A_gW^*(C)A_g^*$, $g\in\Gamma$, form a free family: cyclicity turns any alternating
product into an alternating product of centred elements of $W^*(A)$ and $W^*(C)$. The free group indexes a free family of copies of $S$.
\item \emph{Group labels} (Lemma 2.4). Give a monomial in $A_j^{\pm1}$ and $S$ the product of its group letters as its \emph{label}, ignoring $S$.
Nonidentity labels have trace $0$. Identity labels are polynomials in the $S_g$, obtained by moving letters right with $A_gS=S_gA_g$.
\end{itemize}

\subsection{Move 2: Fox-cocycle flows}
Choose finitely supported real $h_1,\dots,h_n\in\ell^2(\Gamma)$ and extend them to the cocycle $D:\Gamma\to\ell^2(\Gamma)$,
\[
D_{x_j}=h_j,\qquad D_{gb}=D_g+\lambda(g)D_b .
\]
This is Fox's free differential calculus: a word's value is the sum over its letters of their coefficients, shifted by the preceding prefix.
Consider the polynomial ODE, with $S$ held fixed,
\[
A_j'(t)=i\,s_t(h_j)\,A_j(t),\qquad s_t(k)=\sum_gk(g)A_g(t)SA_g(t)^* .
\]
Each letter is rotated on the left by a self-adjoint element built from conjugates of $S$ by its own words. The product rule and the cocycle give
the \emph{word velocity}
\[
A_w'(t)=i\,s_t(D_w)\,A_w(t),
\]
so the cocycle value $D_w$ is the generator of a word's rotation.
\begin{proposition}[the flow; paper's Proposition 3.1]
The ODE has a global unitary solution. It defines a one-parameter group $\beta_t$ of trace-preserving normal automorphisms of $M$ fixing $S$ (hence
$C$). Moreover
\[
\|\beta_t(A_j)-A_j\|\le3\pi|t|\,\|h_j\|_{\ell^2},\qquad \|\beta_t(A_w)-e^{itS}A_w\|\le3\pi|t|\,\|D_w-\delta_e\|_{\ell^2}.
\]
\end{proposition}
\emph{Trace preservation} is the heart. It is proved without knowing freeness at positive times.
\begin{enumerate}[leftmargin=*]
\item On the universal $*$-algebra $\mathcal P$ (unitary $a_j$, self-adjoint $z$), define the derivation $\delta z=0$, $\delta a_j=i\sigma(h_j)a_j$. It respects the
relations, and $\delta z_g=i[\sigma(D_g),z_g]$.
\item \emph{First variation at the free point.} $\delta$ preserves labels, so nonidentity-label terms have trace $0$. On a polynomial $F(z_g)$, the derivative
of $z_g$ is $i[Y_g,S_g]$ with $Y_g=\sum_{b\ne g}D_g(b)S_b$; the own term commutes. Because $S_g$ is free from $B_g=W^*(S_b:b\ne g)\ni Y_g$, conjugating $S_g$ by
$e^{iuY_g}$ preserves the joint law with $B_g$. The first variation is therefore zero, and $\tau\circ\mathrm{ev}_0\circ\delta=0$ on all of $\mathcal P$.
\item \emph{Analytic continuation.} Applied to $\delta^{r-1}p$, step 2 gives $\frac{d^r}{dt^r}\tau(\mathrm{ev}_t p)|_0=\tau(\mathrm{ev}_0\delta^rp)=0$ for every $r\ge1$. Unitarity of the
evolving letters gives the Taylor bound $\frac1{r!}\|\frac{d^r}{dt^r}\mathrm{ev}_tp\|\le b_0R^l(B_1R^J(l+J))^r$, uniformly in $t$. So $t\mapsto\tau(\mathrm{ev}_tp)$ is real-analytic with a
uniform radius, hence constant on $\R$.
\item \emph{Descent.} If $\mathrm{ev}_0p=0$ then $\tau(\mathrm{ev}_t(p^*p))=0$, so the map is well defined. Since $\|b\|=\lim\tau((b^*b)^k)^{1/2k}$ it is isometric. Reversing time
gives surjectivity, and the GNS unitary gives the normal extension.
\end{enumerate}
\emph{Norm control} comes from the free-sum inequality $\|\sum z_j\|\le2(\sum\|z_j\|_2^2)^{1/2}+\max\|z_j\|$ for free centred self-adjoint $z_j$. It gives $\|s(k)\|\le3\pi\|k\|_{\ell^2}$,
\emph{independently of the support size}. Integrating the velocities gives the two bounds, since $s_t(D_w)-S=\beta_t(s(D_w-\delta_e))$.

\subsection{Move 3: the coefficient problem, small letters and a big word}
Take $h_1=h$ and $h_j=0$ for $j\ge2$. Use $a=x_1$, $b_j=x_2^jx_3x_2^{-j}$, $p_j=(ab_1)\cdots(ab_j)$ and $w_m=p_ma$.
\begin{itemize}[leftmargin=*]
\item \emph{Free prefixes} (Lemma 4.2). The $u_j=ab_j$ are free, and $p_j=u_1\cdots u_j$ is a triangular basis change, so $p_1,\dots,p_m$ are free.
\item \emph{Prefix identity.} $D$ vanishes on $\langle x_2,\dots,x_n\rangle$, and $a$ occurs in $w_m$ after the prefixes $p_0,\dots,p_m$. Hence
\[
D_{w_m}=T_mh,\qquad T_m=\sum_{j=0}^m\lambda(p_j).
\]
\item \emph{Spectral limit} (Lemma 4.3). $Q_m=T_mT_m^*$ has $\tau(Q_m^r)=\mathrm{Cat}_r\,(m)_r+O(m^{r-1})$ from non-crossing matchings. So the spectral law of
$Q_m/m$ tends to the free Poisson law $d\nu=\frac1{2\pi}\sqrt{(4-x)/x}\,dx$ on $(0,4)$: $T_m/\sqrt m$ becomes a circular element. Crucially, $\nu(\{0\})=0$.
\item \emph{Cutoff inversion} (Lemma 4.1). $h=T_m^*R_m\delta_e$, with $R_m=Q_m^{-1}\1_{(dm,\infty)}(Q_m)$, gives
\[
T_mh=E_m\delta_e,\qquad \|T_mh-\delta_e\|^2=\mu_m([0,d])\to\nu([0,d])\approx\tfrac2\pi\sqrt d,\qquad \|h\|^2=\tau(R_m)\le\tfrac1{dm}.
\]
Real symmetry and truncation make $h$ real and finitely supported.
\end{itemize}
So a vector of $\ell^2$ norm of order $m^{-1/2}d^{-1/4}$ produces a cocycle value within $\approx d^{1/4}$ of $\delta_e$. The word is $\sim m^2$ letters long.

\subsection{Move 4: one absorption step}
Let $\beta=\beta_1$ for these coefficients, and let $\gamma$ be the automorphism induced by the free-group basis change $(x_1,\dots,x_n,c)\mapsto(x_1,\dots,x_n,cw)$. Then
$\gamma(A_j)=A_j$ and $\gamma(C)=CA_w$. Put $\alpha=\gamma^{-1}\circ\beta$. Exactly,
\[
\alpha(A_j)-A_j=\gamma^{-1}(\beta(A_j)-A_j),\qquad \alpha(A_w)-C=\gamma^{-1}(\beta(A_w)-CA_w),
\]
so $\max_j\|\alpha(A_j)-A_j\|<\varepsilon$ and $\|\alpha(A_w)-C\|<\varepsilon$. The new $n$-tuple is $\varepsilon$-close to the old one, and one of its words is $\varepsilon$-close to the old $C$.
The new completion is $C'=CA_w^*$.

\subsection{Move 5: the adaptive limit}
Fix a dense sequence $(y_j)$ in $L^2(M)$. At stage $k$ the choices are made in this order:
\begin{enumerate}[leftmargin=*]
\item approximate $y_1,\dots,y_k$ to $2^{-k}$ by polynomials $F_{j,k}$ in the current $(n+1)$-tuple;
\item choose $\varepsilon_k<\min(r_{k-1}/2,\,2^{-k}/(1+H_k))$, with $H_k$ the $L^2$-Lipschitz constant of the $F_{j,k}$;
\item absorb (Move 4), obtaining the word $w_k$;
\item substitute $G_{j,k}(X)=F_{j,k}(X,X_{w_k})$, which are polynomials in the $n$-tuple alone;
\item only now choose the future budget $r_k\le\min(r_{k-1}/2,\,2^{-k}/(1+L_k))$, where $L_k$ reflects the length of $w_k$.
\end{enumerate}
Then $\sum_{\ell>k}\varepsilon_\ell\le r_k$, so the $n$-tuples converge in operator norm to $A^{(\infty)}$. Each word's trace is norm-continuous, so $A^{(\infty)}$ is free Haar. Moreover
$\|G_{j,k}(A^{(\infty)})-y_j\|_2\le3\cdot2^{-k}$, so the word algebra of $A^{(\infty)}$ is $L^2$-dense. The trace-preserving conditional expectation onto $W^*(A^{(\infty)})$ then fixes a dense set,
so $W^*(A^{(\infty)})=M$. The completions $C^{(k)}$ need not converge.

\section{The essence: four mechanisms}
\label{sec:essence}
\subsection{Mechanism 1: a gauge field on words, and why it must be outer}
The flow rotates each letter on the left by a self-adjoint element built from conjugates of $S$. A word then rotates by the
\emph{accumulated} element $s(D_w)$, the Fox cocycle evaluated on the word. This is a gauge field along words: $D$ plays the connection, $D_w$ its
holonomy, and $e^{iS}=C$ is the target holonomy.

Inner automorphisms fixing $C$ cannot do this. $W^*(C)$ is a maximal abelian subalgebra of $M$, so they are $\Ad u$ with $u=e^{if(S)}$, and they
send $A_j\mapsto e^{if(S)}e^{-if(S_{x_j})}A_j$. Those are coboundary coefficients, $D_g=(1-\lambda(g))v$ up to the functional calculus. A coboundary with
$D_{x_j}=0$ for $j\ge2$ forces $v$ to be invariant under $x_2,\dots,x_n$, hence $v=0$. The absorption therefore uses \emph{non-trivial} cocycles in
$H^1(\F_n;\ell^2\F_n)$.

That group is non-zero because $\beta^{(2)}_1(\F_n)=n-1>0$. There is an irony here. The first $L^2$-Betti number, long hoped to be the invariant that
\emph{distinguishes} the free group factors, is exactly what supplies the cocycles that make them isomorphic.

\subsection{Mechanism 2: coherent $\sqrt m$ amplification along free prefixes}
For $h$ concentrated near $e$ and free prefixes far apart, $\|T_mh\|^2\approx(m+1)\|h\|^2$. Prefix shifts are nearly orthogonal, so the word velocity
grows like $\sqrt m$ times the letter velocity even without design: this is a free central limit, $T_m/\sqrt m\to$ circular. Design does the rest. The
minimal-norm solution $h=T_m^*(T_mT_m^*)^{-1}\delta_e$, with the inverse truncated, hits $\delta_e$ at cost $\|h\|^2=\tau(R_m)\approx2/(\pi m\sqrt d)$.

This works only because the free Poisson law has \emph{no atom at zero}, so $T_m^*$ has no kernel in the direction of $\delta_e$. Letters are moved by
$O(m^{-1/2}d^{-1/4})$ while a word of length $\sim m^2$ is rotated by $C$ up to $O(d^{1/4})$.

\subsection{Mechanism 3: freeness turns $\ell^2$ into operator norm}
For bounded \emph{commuting} independent variables, $\|\sum k_gX_g\|_\infty=\sum|k_g|\,\|X_g\|_\infty$ is an $\ell^1$ quantity, since all extremes can coincide. For free
variables, $\|\sum k_gS_g\|\le2\|S\|_2\|k\|_{\ell^2}+\|S\|\max|k|$, an $\ell^2$ quantity: free variables cannot conspire. This is what lets a coefficient vector of huge
support but tiny $\ell^2$ norm move every letter by a tiny amount in operator norm. In a commutative setting the same $h$ would move letters by
$\sim\sqrt{\text{support}}\cdot\|h\|_{\ell^2}$, which is large.

\emph{Non-commutativity is the resource.} The free family $\{S_g\}$ indexed by the free group is what makes small-but-spread vector fields
small.

\subsection{Mechanism 4: distribution is closed, generation is not}
Having free Haar word distribution is a set of trace conditions on words, each continuous in operator norm, so it survives norm limits.
Generating $M$ is a density condition, which is neither open nor closed under norm limits. Each $A^{(k)}$ generates a proper subfactor with a
free complement, yet the limit generates everything.

The adaptive budgets are what make generation survive. Each target $y_j$, once approximated by a polynomial in the current $n$-tuple, is
protected by all later budgets, which are chosen only after the word length that the polynomial needs is known. It is a diagonal argument in
which the complexity of the witnesses (word lengths) grows without bound while the moves shrink faster.

\subsection{The whole argument in one line}
\begin{multline*}
\underbrace{\beta^{(2)}_1(\F_n)>0}_{\text{outer Fox flows}}\ +\ \underbrace{\nu_{\rm FP}(\{0\})=0}_{\text{small letters, big word}}\ +\ \underbrace{\text{free Khintchine}}_{\ell^2\to\|\cdot\|}\\
+\ \underbrace{\text{adaptive budgets}}_{\text{generation in the limit}}\ \Longrightarrow\ L(\F_n)\cong L(\F_{n+1}).
\end{multline*}

\section{Why it evaded thirty years, and what changes}
\label{sec:history}
\subsection{The direction that failed}
For thirty years the main line of attack sought an \emph{invariant that counts generators}. Candidates were Voiculescu's free entropy
dimensions ($\delta,\delta_0,\delta^*,\delta^\star$), Connes--Shlyakhtenko $L^2$-Betti numbers of the algebra, and cost-like quantities. On canonical generators they
all return $n$ (or $n-1$). The expert's remark, that ``we have been heading in the wrong direction'', reads naturally against this background. The paper
shows that the count is a property of the \emph{generating set}, not of the algebra:
\begin{itemize}[leftmargin=*]
\item a free Haar $n$-tuple, with $\delta_0=n$, generates $L(\F_{n+1})$;
\item $L(\F_2)$ has finite generating tuples with $\delta=\delta_0=\delta^*=\delta^\star=n$ for every $n\ge2$ (Corollary 7.1).
\end{itemize}
The resolution came instead from \emph{dynamics}: an explicit, rich family of automorphisms and an adaptive limit. That is a constructive object, not a
numerical invariant.

\subsection{What stays true, and what must change}
\begin{itemize}[leftmargin=*]
\item \emph{Upper-bound theorems survive.} These are results that bound $\delta_0$ for \emph{all} generating sets, such as Voiculescu's absence of Cartan
subalgebras, Ge's primeness, Jung's strong $1$-boundedness, and Hayes' $1$-bounded entropy. They are statements about every generating set, and remain
valid. $L(\F_n)$ still has no Cartan subalgebra, is prime and strongly solid, and $h(L(\F_n))=\infty$ for all $n$.
\item \emph{Generator-counting invariants must collapse.} Any quantity that is an invariant of the von Neumann algebra and was expected to equal
$n$ or $n-1$ on $L(\F_n)$ must take a single value for all $n$. For example, the Connes--Shlyakhtenko first $L^2$-Betti number of $L(\F_n)$ cannot
equal $n-1$ for every $n$.
\item \emph{Fundamental group.} $\mathcal F(L(\F_2))=\R_{>0}$; every interpolated factor is isomorphic to every other, including $L(\F_\infty)$.
\item \emph{Generation is not stable under small perturbations.} Choosing $r_0$ small in the construction gives free Haar $n$-tuples \emph{arbitrarily close in
operator norm} to the canonical generators $\lambda(x_1),\dots,\lambda(x_n)$ of $L(\F_{n+1})$ that generate all of it. The canonical tuple generates a proper
subfactor with a free complement. Generation is invisible to any norm-local count.
\end{itemize}

\subsection{Relation to earlier techniques}
\begin{itemize}[leftmargin=*]
\item \emph{Voiculescu's cyclomorphy} exponentiates trace-preserving polynomial vector fields on free semicircular generators. The paper redoes it for
Haar unitaries with a fixed bounded logarithm, proving descent through the unitary relations from moments.
\item \emph{Guionnet--Shlyakhtenko free monotone transport} builds isomorphisms by changes of variables at a \emph{fixed} rank. Here the rank changes
through an infinite product of small automorphisms and a limiting generation argument.
\item \emph{Shlyakhtenko's Fuchsian preprint} (September 2026) finds free complements by finite automorphisms plus limiting generation. The idea was
in the air, and the small-cocycle absorption is the new ingredient.
\item \emph{Fox calculus} (1953) supplies the cocycle rule, and \emph{Ricard--Xu} the free Khintchine decomposition.
\end{itemize}

\section{Critical reading}
\label{sec:critical}
We went through every step, looking for the places where a proof of this kind usually breaks.
\begin{center}\small
\begin{tabular}{>{\raggedright\arraybackslash}p{4.4cm}>{\raggedright\arraybackslash}p{9.6cm}}\toprule
step & what could fail, and why it does not\\\midrule
Haar-tuple identification (L.\,2.2) & needs only orthonormality of words and density of their span; both follow from the trace conditions and generation\\
free conjugates (L.\,2.3) & cyclic rotation turns $\tau(\prod A_{g_j}c_jA_{g_j}^*)$ into an alternating centred product of free algebras; correct\\
labels (L.\,2.4) & the alternating-reduction induction is standard; nonidentity labels survive the removal of scalar parts\\
derivation on $\mathcal P$ & respects $a_j^*a_j=1$ since $\sigma(h_j)$ is self-adjoint ($h_j$ real); the cocycle rule follows from the product rule\\
first variation (L.\,3.2) & the own term commutes; $Y_g\in B_g$; conjugating one free variable by a unitary of its free complement preserves the joint law. Freeness is
used \emph{only at $t=0$}\\
higher derivatives & $\tau\circ\mathrm{ev}_0\circ\delta=0$ holds on all of $\mathcal P$, hence on $\delta^{r-1}p$\\
analyticity (L.\,3.4) & the Taylor bound is uniform in $t$ because the evolved letters are unitary and $S$ is fixed; constancy on $\R$ follows from real-analytic
continuation\\
descent, isometry, surjectivity & faithfulness of $\tau$; $\|b\|=\lim\tau((b^*b)^k)^{1/2k}$; uniqueness of the reversed flow\\
norm bound (L.\,2.5) & creation, annihilation and diagonal decomposition on the free product $L^2$; constant $3\pi$ (we measure $2\pi/\sqrt3$ at the free CLT edge)\\
prefix identity (L.\,4.2) & checked exactly by computer (Section~\ref{sec:checks})\\
free Poisson limit (L.\,4.3) & non-crossing matching count plus tightness and moment determinacy on $[0,4]$; we observe it numerically\\
cutoff inversion (L.\,4.1) & order of choices $d\to m\to$ finite truncation; real symmetry via conjugation commuting with real functional calculus\\
absorption (P.\,5.1) & exact algebra with $\gamma(C)=CA_w$; the error transfers isometrically\\
adaptive limit (Thm 1.1) & $\varepsilon_k$ chosen before $w_k$, $r_k$ after; $\sum_{\ell>k}\varepsilon_\ell\le r_k$; $L^2$-density plus conditional expectation gives generation\\\bottomrule
\end{tabular}
\end{center}
Further remarks.
\begin{itemize}[leftmargin=*]
\item \emph{Rates.} With the actual spectral asymptotics, $\|h\|^2\approx2/(\pi m\sqrt d)$ and error$^2\approx(2/\pi)\sqrt d$. Accuracy $\eta$ needs $d\sim(\pi\eta^2/2)^2$ and
$m\sim4/(\pi^2\eta^4)$, with word length $\sim m^2\sim\eta^{-8}$. That is polynomial blow-up, irrelevant for existence and a mild cost for any numerical
realization.
\item \emph{Constructivity.} Each stage is explicit, with an explicit word, coefficient vector and flow. The limit is an infinite product, and the
isomorphism has no closed form.
\item \emph{External inputs.} Only Dykema's amplification formula is used for Theorem~1.2, and only the alternative for Corollary~6.1. Both are classical.
\item \emph{What a referee will read hardest.} The first-variation argument, and the claim that real-analytic continuation needs no freeness
at $t>0$. We find both correct. The authors' care in separating formal identities on $\mathcal P$ from evaluated ones is what makes the second step
legitimate.
\end{itemize}

\section{Checks}
\label{sec:checks}
The script \texttt{scripts/\allowbreak verify\_\allowbreak ffac.py} runs in a few minutes. Its transcript is
\texttt{notes/\allowbreak stage17/\allowbreak verify\_\allowbreak ffac.txt}.
\begin{center}\small
\begin{tabular}{>{\raggedright\arraybackslash}p{5.0cm}>{\raggedright\arraybackslash}p{9.0cm}}\toprule
statement & result\\\midrule
prefix identity $D_{w_m}=T_mh$ (L.\,4.2) & exact (difference $0$) for $m=1,3,5$, with word lengths $5,19,41$\\
cocycle identity $D_{uv}=D_u+\lambda(u)D_v$ & exact on $200/200$ random pairs\\
free prefixes & $936$ reduced words of length $\le4$ in $p_1,p_2,p_3$: all distinct, none trivial\\
free Poisson limit (L.\,4.3), $N=1200$ & moments of $Q_m/m$ at $m=128$: $1.007,2.021,5.060,14.18$ (Catalan $1,2,5,14$); $\mu([0,0.05])=0.1417$ vs $\nu=0.1421$; top eigenvalue $3.95$ (edge $4$)\\
cutoff inversion (L.\,4.1) & $m=128$, $d=0.02$: $\|h\|^2=0.031$ (bound $0.39$), error$^2=0.089$ ($\nu([0,d])=0.090$)\\
rates & $\eta=0.3$: $m\approx50$; $\eta=0.1$: $m\approx4\times10^3$, word length $\approx1.6\times10^7$\\
free-sum bound (L.\,2.5), $N=600$ & spread coefficients: $20.1$ against bound $52.6$ and free-CLT edge $20.2$; one large coefficient: $30.6$ against $81.9$\\
flow mechanism (P.\,3.1), $N=160$ & $m=16$: letter velocity $0.93$, word deviation $\|s(D_w)-S\|=1.38$; $m=64$: $0.74$ and $1.02$; $\|S\|=3.13$; within the $3\pi$ bounds\\
amplification arithmetic & $N_3^{\sqrt2}=N_2$, $N_5^{\sqrt2}=N_3$\\\bottomrule
\end{tabular}
\end{center}
In a matrix model, conjugation fixes the identity matrix, so $S$ must be centred and the trivial direction removed. In the free setting the
conjugates of the centred $S$ span a copy of $\ell^2(\Gamma)$, which has no such direction. The mechanism check shows the paper's trade-off: as $m$
grows, letters move less while the word's rotation approaches the pure rotation by $S$. Convergence is slow, as the rates predict.

\section{What it means for our programme}
\label{sec:programme}
\subsection{The same cocycle}
The paper's Fox cocycle and our stage-15 error cocycle are one structure:
\[
D_{gb}=D_g+\lambda(g)D_b\qquad\longleftrightarrow\qquad \delta_l=T_l\delta_{l-1}+\beta_l,\quad \delta_L=\sum_lT_L\cdots T_{l+1}\beta_l .
\]
In both, a word's value is the sum of local contributions, each transported by part of the word: the prefix (left regular action) in theirs, the
downstream gated maps in ours. Their analytic heart concerns the Gram of the transports. For free prefixes it tends to a circular element,
with no atom at zero. Ours is what the measured sum rule reports: births from different layers add in quadrature, with transport factors $0.65$--$0.94$.
\begin{itemize}[leftmargin=*]
\item \emph{Incoherent local contributions add in $\ell^2$.} The word grows like $\sqrt m$. This is our late, incoherent residual.
\item \emph{Designed contributions can be coherent.} The minimal-norm input $T^*(TT^*)^{-1}\cdot\mathrm{target}$ reaches an $O(1)$ target with local inputs of size $m^{-1/2}$.
For an estimator this is the adjoint (dual-weighted) design of corrections, and it is efficient exactly when the transport Gram has no small
singular values in the target direction.
\end{itemize}

\subsection{The same orthogonality}
Free Khintchine, $\|\sum k_gS_g\|\lesssim\|k\|_{\ell^2}$, is the non-commutative form of our row-Mehler theorem: independence across rows makes errors add in
$\ell^2$ rather than $\ell^1$. In the paper this is a \emph{resource}, since spread corrections are small. For us it is a \emph{constraint}, since spread errors
cannot be cancelled by a few global corrections (stage~15, Corollary~3.3).

\subsection{The same lesson about counting}
The paper shows that a generator count is a property of the generating set, not of the algebra; the invariant structure is dynamical. Our
stage~16 measured the analogue for estimators. A representation in a fixed basis needs degree $\sim17$ at depth~16, and the tractable
representation co-evolves with the weights. In both cases the productive object is a \emph{flow of frames} adapted to the instance, not a count.

\subsection{A template for dynamic interdependence}
The paper's iteration is the cleanest rigorous example we know of a computation whose next step is fixed only after the previous output
is seen. Each pass makes a small, structure-preserving move that serves a target, and the tolerance of every future move is set by the
sensitivity of everything already achieved. The transferable design principle for an estimator has three parts:
\begin{itemize}[leftmargin=*]
\item \emph{Passes of small, exact moves.} A sequence of passes, each a small move of the internal frame that preserves the exact identities
(homogeneity, Euler--Stein, moment positivity). This is the analogue of trace preservation.
\item \emph{Moves aimed through the adjoint.} Each move is directed at the currently unresolved part, the analogue of absorbing $C$.
\item \emph{Budgets set by sensitivity.} Each pass's tolerance comes from the adjoint sensitivities of what earlier passes achieved, the analogue
of $r_k$.
\end{itemize}

\emph{What does not transfer.} Mechanism~3 is non-commutative. Our network's state lives in a commutative probability space, where spread moves
cost $\ell^1$, so the absorption phenomenon itself has no classical analogue. What transfers is the cocycle algebra, the spectral (Gram)
analysis of transports, and the order of choices in the adaptive iteration.

\section{Addendum: deductions from the companion proof study}
\label{sec:deductions}
The companion proof study (\emph{Absorbing a free generator}, 13 pp) reconstructs the argument independently. It adds three deductions
(its 8.1--8.3) and reads $M$ as a crossed product (its \S9). We confirm all four. We also sharpen three of them and add one quantitative law:
\begin{itemize}[leftmargin=*]
\item a relative form of the theorem that fixes a subfactor pointwise;
\item weak nullness of the completions;
\item a hard-edge steering law that fixes the cost exponent of the paper's route;
\item the identification of the Fox flows as cocycle self-conjugacies of the free Bernoulli shift.
\end{itemize}
Two verification caveats from the study stand. We have not rerun the reported Lean formalization, and the quoted expert is not
authenticated. Our verdict rests on our own reading.

\subsection{Only one letter moves (study 8.1), and a relative absorption}
In Move~3, $h_1=h$ and $h_j=0$ for $j\ge2$, so $A_j'=0$ for $j\ge2$ and $\beta$ fixes $A_2,\dots,A_n$. The basis change $\gamma$ fixes every $A_j$, so $\alpha$
fixes $A_2,\dots,A_n$. Every stage of the adaptive limit does the same with the current first letter. The limit tuple is therefore
$(B_1,A_2,\dots,A_n)$ with $\|B_1-A_1\|<\delta$.
\begin{corollary}[relative absorption]
Let $n\ge3$, $P=W^*(A_2,\dots,A_n)\cong L(\F_{n-1})$ and $k\ge1$. There is an isomorphism $P*L(\Z)\cong P*L(\F_{k+1})$ that is the identity on $P$ and
sends the generator of $L(\Z)$ to within any prescribed $\delta$ of $A_1$. In particular a fixed copy of $L(\F_{n-1})$ has free complements
$W^*(B_1)\cong L(\Z)$ (abelian) and $W^*(A_1,C)\cong L(\F_2)$ (a factor) in the same algebra. Free complements are not unique even up to isomorphism.
\end{corollary}
\begin{proof}
For $k=1$, apply the theorem to the $n$-tuple $(A_1,\dots,A_n)$ with completion $C$. Only $A_1$ moves, and the limit tuple is free Haar, so $B_1$ is
free from $P$. For general $k$, start from $(A_1,\dots,A_n,C_1,\dots,C_k)$ and absorb $C_1,\dots,C_k$ in turn into the current first letter, using the
tuple $(B_1^{(i)},A_2,\dots,A_n,C_{i+2},\dots,C_k)$ at step $i+1$. The prefix words use only the first three letters, so $n\ge3$ suffices, and the
$\delta_i$ can be summed.
\end{proof}

\subsection{Strong but never uniform recovery (study 8.2), and an escaping completion}
Let $\cA_k=W^*(A^{(k)})$ and let $E_k$ be its trace-preserving conditional expectation. The study proves $E_k\to I$ strongly on $L^2(M)$:
\begin{itemize}[leftmargin=*]
\item $\|E_ky_j-y_j\|_2\le\|G_{j,k}(A^{(k)})-y_j\|_2\to0$ on a dense set;
\item $\|E_k\|\le1$ extends this to all of $L^2(M)$.
\end{itemize}
It also proves $\|I-E_k\|=1$: $C^{(k)}$ is a Haar unitary free from $\cA_k$, so $E_kC^{(k)}=\tau(C^{(k)})=0$ with $\|C^{(k)}\|_2=1$. We add:
\begin{proposition}
$C^{(k)}\to0$ weakly in $L^2(M)$. Hence the completions have no $L^2$ limit, where the paper's ``need not converge'' leaves this open. In the
ultrapower, $(C^{(k)})_\omega\in M^\omega$ is a Haar unitary free from $M$.
\end{proposition}
\begin{proof}
For $y\in L^2(M)$, $|\langle C^{(k)},y\rangle|=|\langle C^{(k)},(I-E_k)y\rangle|\le\|(I-E_k)y\|_2\to0$.

For freeness, take an alternating product of centred $y_i\in M$ and centred polynomials in $C^{(k)}$. Replace each $y_i$ by the centred, norm-bounded
$E_k(y_i)$. This changes the trace by at most $\sum_i\|y_i-E_ky_i\|_2\prod_{j\ne i}\|x_j\|\to0$, and the replaced product has trace $0$ by freeness at stage
$k$.
\end{proof}
Popa showed that every separable II$_1$ factor has such a free Haar unitary in $M^\omega$. Here the construction produces one explicitly, as the
discarded completions. Recovery is target by target: every fixed observable is recovered, while a unit direction, the current completion, is
always invisible and drifts off to infinity.

\subsection{Word length (study 8.3), refined, and a steering law}
$C$ is free from $W^*(A)$, and $A_w$ is a Haar unitary for $w\ne e$. So $C^*A_w$ is Haar, its spectrum is the whole circle, and $\|C-A_w\|=2$. The study
telescopes this against $\varepsilon$-moves of the letters to get $|w|>2/\varepsilon-1$. By 8.1 only occurrences of $x_1^{\pm1}$ cost anything.
\begin{proposition}[universal bound]
If an absorption moves only the first letter, by less than $\varepsilon$, and $\|B_w-C\|<\varepsilon$, then $\ell_1(w)>2/\varepsilon-1$, where $\ell_1(w)$ is the number of
occurrences of $x_1^{\pm1}$ in $w$.
\end{proposition}
In the construction $\ell_1(w_m)=m+1$, while $|w_m|=m^2+3m+1$: the letters $b_j=x_2^jx_3x_2^{-j}$ are padding that keeps the prefixes free. How
large must $m$ be on the paper's route? The answer is set by the coefficient problem of Move~3.

\begin{proposition}[steering law]
Let $\mu_m$ be the spectral measure of $Q_m=T_mT_m^*$ at $\delta_e$. For every $h$ and every $\lambda>0$,
\[
\|T_mh-\delta_e\|^2+\lambda\|h\|^2\ \ge\ \int\frac{\lambda}{q+\lambda}\,d\mu_m(q),
\]
with equality at the Tikhonov solution $h=T_m^*(Q_m+\lambda)^{-1}\delta_e$. Take $\lambda=md$ with $1/m\ll d\ll1$ and replace $\mu_m(m\,\cdot)$ by its free Poisson limit
$\nu$, whose density is $\sim1/(\pi\sqrt x)$ near $0$. The right side is then $\sqrt d\,(1+o(1))$. Optimising over $d$ gives the Pareto frontier
\[
\|T_mh-\delta_e\|_{\ell^2}\cdot\|h\|_{\ell^2}\ \ge\ \frac{1+o(1)}{2\sqrt m}.
\]
Tikhonov attains it. The paper's spectral cutoff attains $2/(\pi\sqrt m)$, within a factor $4/\pi$.
\end{proposition}
\begin{proof}
The inequality is the minimum of a quadratic, with value $\lambda\langle\delta_e,(Q_m+\lambda)^{-1}\delta_e\rangle$. Also
$\int_0^\infty\frac{d}{x+d}\frac{dx}{\pi\sqrt x}=\sqrt d$. With $a=m\|h\|^2$, $\|T_mh-\delta_e\|^2\ge\sup_d(\sqrt d-da)=1/(4a)$.

For the cutoff, $\nu([0,d])\approx\frac2\pi\sqrt d$ and $\int_d\frac{d\nu}x\approx\frac2{\pi\sqrt d}$, whose product is $4/\pi^2$.
\end{proof}
\begin{corollary}[cost of the Fox-prefix route]
Asking for letter moves and the word error both below $\varepsilon$, with the proven constant $3\pi$, is sufficient, to leading order, once
$m\ge81\pi^4/(4\varepsilon^4)$. To first order, $\|s(k)\|\ge\|S\|_2\|k\|_2=(\pi/\sqrt3)\|k\|_2$, so $m\gtrsim\pi^4/(36\varepsilon^4)$ is necessary on this route.

Hence the route needs $m=\Theta(\varepsilon^{-4})$ moved letters and $|w|=\Theta(\varepsilon^{-8})$, against the universal $\ell_1(w)>2/\varepsilon-1$.
\end{corollary}
\begin{remark}[the exponent is the hard edge]
Suppose the Gram law near $0$ is $\nu([0,x])\sim c\,x^\alpha$ with $0<\alpha<1$. The same Legendre argument gives
$m\gtrsim\|h\|^{-2}\,\mathrm{err}^{-2(1-\alpha)/\alpha}$.
\begin{itemize}[leftmargin=*]
\item A spectral gap costs nothing in accuracy.
\item An atom at $0$ makes the target unreachable.
\item The free Poisson law has the square-root hard edge $\alpha=\frac12$, which is what fixes the paper's $\varepsilon^{-4}$.
\end{itemize}
With free prefixes, shorter padding could shorten $|w|$ but cannot change the exponent in $m$. Whether a different absorption mechanism closes the gap between $\varepsilon^{-1}$ and $\varepsilon^{-4}$ is open.
\end{remark}

\subsection{The crossed product (study \S9), and what the flows are in it}
Let $\cN=W^*(S_g:g\in\Gamma)$. By Lemma~2.3 it is the free product $*_{g\in\Gamma}W^*(S_g)$ of copies of $L(\Z)$. The unitaries $A_g$ normalise it,
$\sigma_g:=\Ad A_g$ shifts the index, and by Lemma~2.4 $E_\cN(A_g)=0$ for $g\ne e$. Together with $C=e^{iS}\in\cN$, this gives the study's
\[
M=\cN\rtimes_\sigma\Gamma\qquad\text{(the free Bernoulli shift of $\Gamma=\F_n$ with base $L(\Z)$)},
\]
and $L^2(M)=\bigoplus_g L^2(\cN)A_g$, graded by labels. The theorem then reads $(*_\Gamma L(\Z))\rtimes\Gamma\cong L(\Gamma)$: the crossed product absorbs its
coefficient algebra into the group algebra.
\begin{proposition}[Fox flows are cocycle self-conjugacies]
For all $t$, $\beta_t(\cN)=\cN$ and $\beta_t(A_g)=U_g(t)A_g$ with $U_g(t)\in\mathcal U(\cN)$. Moreover:
\begin{itemize}[leftmargin=*]
\item $U_{gb}=U_g\,\sigma_g(U_b)$, so $U(t)$ is a unitary $1$-cocycle of the shift;
\item $\theta_t=\beta_t|_\cN$ satisfies $\theta_t\sigma_g\theta_t^{-1}=\Ad U_g(t)\,\sigma_g$ and $\theta_t(S)=S$;
\item $\frac d{dt}U_g(t)|_{0}=i\,s(D_g)$.
\end{itemize}
So $\beta_t$ preserves the grading. The Fox cocycle, pushed through the equivariant map $s:\ell^2\Gamma\to\cN$ (with $\sigma_g\circ s=s\circ\lambda(g)$), is the tangent of a
unitary cocycle of the free Bernoulli shift.
\end{proposition}
\begin{proof}
The derivation $\delta$ preserves labels. So $\mathrm{ev}_0(\delta^ra_j)$ lies in $\cN A_j$, because a label-$x_j$ monomial times $A_j^*$ is a polynomial in the $S_g$ by
Lemma~2.4. The Taylor series converges in norm for $|t|<R$, with the uniform radius of Lemma~3.4. Since $\cN A_j$ is norm closed,
$\beta_t(A_j)\in\cN A_j$ for $|t|<R$.

The cocycle rule follows from $\beta(A_gA_b)=U_gA_gU_bA_b$. Then $\beta_t(S_g)=U_gS_gU_g^*\in\cN$, and with $\beta_{-t}$ this gives $\beta_t(\cN)=\cN$. The group law
extends everything to all $t$. The derivative is the word velocity of Move~2 at $t=0$.
\end{proof}
The basis change $\gamma$ is the one move that breaks the grading. $\gamma(C)=CA_w$ lies in the $w$-graded piece $\cN A_w$, so $\gamma(\cN)\ne\cN$. Each absorption
step is therefore:
\begin{enumerate}[leftmargin=*]
\item a continuous, grading-preserving cocycle move $\beta$, which steers the $w$-component of $A_w$ to $\approx e^{iS}A_w=CA_w$;
\item a discrete regrading $\gamma^{-1}$ by $w$.
\end{enumerate}
The rank changes only through the discrete regradings, and each one costs $\ell_1(w)>2/\varepsilon-1$.

We agree with the study's caveats. $\beta_t$ is not a modular flow, since $\tau$ is a trace and the modular group is trivial. $D$ is a group cocycle,
not a Dirac operator. An abstract isomorphism forgets the marked generators, and with them all geometric and computational data.

\subsection{For the programme}
\begin{itemize}[leftmargin=*]
\item \emph{A steering law for our transports (conjecture, to measure).} The stage-15 no-free-lunch for local corrections should have the quantitative
form above. Cancelling a target at depth gap $k$ with corrections of size $\|h\|$ should cost $\|h\|^{-2}\mathrm{err}^{-2(1-\alpha_k)/\alpha_k}$, where $\alpha_k$ is the
hard-edge exponent of the gated transport Gram $T_L\cdots T_{L-k+1}(\cdot)^*$ in the output-mean direction. Products of $k$ free circular factors have
Fuss--Catalan edges $\alpha=1/(k+1)$, which would predict a cost $\mathrm{err}^{-2k}$. Measuring $\alpha_k$ on net~0 from the stage-16 Jacobians is a single
quick check.
\item \emph{The right notion of convergence for an instance-coupled estimator} is 8.2's: strong, target by target, never uniform. A frame fixed
before the target is known has a blind unit direction. The adaptive order of choices is what turns strong convergence into recovery.
\item \emph{Continuous frame motion does not change rank.} The co-evolving frames of stage~16 are the analogue of the grading-preserving
cocycle moves. Changing \emph{what} is carried requires a separate, discrete regrading move. Its decoder complexity grows at least like
$1/\varepsilon$, and on linear-response routes like the hard-edge law. This is a reading, not a theorem, for our commutative setting.
\end{itemize}

\end{document}
